%% ============================================================
%% shared/well_architected.tex
%% The AWS Well-Architected Framework — underpins BOTH exams
%% ============================================================

\servicesection{AWS Well-Architected Framework}{\bothtag}{%
  The six pillars that guide every SAA and SAP architecture question.}

The Well-Architected Framework gives you a consistent way to evaluate architectures
and implement designs that scale over time. \textbf{Every exam domain maps back to
at least one pillar.}

\subsection{The Six Pillars}

\begin{keyfacts}
\begin{enumerate}
  \item \textbf{Operational Excellence} — run, monitor, and improve systems.
        Focus on automating operations and responding to events.
  \item \textbf{Security} — protect data, systems, and assets. Includes IAM,
        detective controls, infrastructure protection, and incident response.
  \item \textbf{Reliability} — recover from failures and meet demand dynamically.
        Covers HA, DR, and fault tolerance.
  \item \textbf{Performance Efficiency} — use resources efficiently as demand
        changes; choose the right resource type and size.
  \item \textbf{Cost Optimization} — avoid unnecessary costs; match supply to
        demand; use the right pricing model.
  \item \textbf{Sustainability} — minimise environmental impact of cloud workloads.
        (Newer pillar — appears more on SAP.)
\end{enumerate}
\end{keyfacts}

\begin{examtip}
When an exam question asks you to choose between two architecturally valid options,
ask yourself: \emph{which pillar is the question testing?} Keywords like
"least operational overhead" → Cost/Operational Excellence;
"most fault-tolerant" → Reliability; "most secure" → Security.
\end{examtip}

\subsection{Design Principles (per pillar)}

\subsubsection{Operational Excellence}
\begin{itemize}
  \item Perform operations as code (Infrastructure as Code)
  \item Make frequent, small, reversible changes
  \item Anticipate failure (game days, chaos engineering)
  \item Learn from all operational events and failures
\end{itemize}

\subsubsection{Security}
\begin{itemize}
  \item Implement a strong identity foundation (least privilege)
  \item Enable traceability (CloudTrail, Config, GuardDuty)
  \item Apply security at all layers
  \item Protect data in transit and at rest (KMS, TLS)
  \item Keep people away from data (automate where possible)
\end{itemize}

\subsubsection{Reliability}
\begin{itemize}
  \item Test recovery procedures
  \item Automatically recover from failure
  \item Scale horizontally to increase aggregate system availability
  \item Stop guessing capacity — use Auto Scaling
  \item Manage change using automation
\end{itemize}

\subsubsection{Performance Efficiency}
\begin{itemize}
  \item Democratise advanced technologies (use managed services)
  \item Go global in minutes (multi-region)
  \item Use serverless architectures
  \item Experiment more often
  \item Consider mechanical sympathy (use the right tool for the job)
\end{itemize}

\subsubsection{Cost Optimization}
\begin{itemize}
  \item Implement cloud financial management (FinOps)
  \item Adopt a consumption model (pay for what you use)
  \item Measure overall efficiency (CloudWatch metrics per \$)
  \item Stop spending money on undifferentiated heavy lifting
  \item Analyse and attribute expenditure (tagging, Cost Explorer)
\end{itemize}

\subsubsection{Sustainability}
\begin{itemize}
  \item Understand your impact (carbon footprint tool)
  \item Maximise utilisation (right-size, use Graviton)
  \item Adopt managed services (AWS manages hardware lifecycle)
\end{itemize}

\begin{sapdepth}
At the SAP level, questions tie the framework to \textbf{organisational strategy}:
how do you enforce pillars across multiple accounts? Key tools:
\begin{itemize}
  \item \awsservice{AWS Well-Architected Tool} — workload reviews, milestone tracking
  \item \awsservice{AWS Trusted Advisor} — automated best-practice checks (5 categories map to pillars)
  \item \awsservice{AWS Control Tower} — landing zone that enforces guardrails (SCPs) at scale
  \item \awsservice{AWS Service Catalog} — approved product portfolios → governance at scale
\end{itemize}
\end{sapdepth}

\begin{examtip}[title={\faLightbulb~SAP exam tip: "most operationally efficient"}]
SAP questions often say "with the \textbf{least operational overhead} going forward."
This almost always points to a \textbf{managed/serverless} service over a self-managed
EC2 solution. E.g.\ Aurora Serverless over RDS on EC2, or Fargate over EC2-backed ECS.
\end{examtip}
